% Generated draft: requires mathematical and source review.
\documentclass[11pt]{article}
\usepackage[margin=25mm]{geometry}
\usepackage{fontspec}
\setmainfont{Latin Modern Roman}
\usepackage{amsmath,amssymb,mathtools,unicode-math}
\setmathfont{Latin Modern Math}
\usepackage{enumitem,needspace,xurl,hyperref}
\setlength{\parindent}{0pt}
\setlength{\parskip}{5pt}
\setlength{\emergencystretch}{3em}
\newcommand{\N}{\mathbb{N}}
\newcommand{\Z}{\mathbb{Z}}
\newcommand{\Q}{\mathbb{Q}}
\newcommand{\R}{\mathbb{R}}
\newcommand{\C}{\mathbb{C}}
\newcommand{\F}{\mathbb{F}}
\newcommand{\A}{\mathbb{A}}
\newcommand{\PP}{\mathbb{P}}
\newcommand{\E}{\mathbb{E}}
\newcommand{\eps}{\varepsilon}
\newcommand{\dd}{\,\mathrm{d}}
\begin{document}
\section*{UnsolvedMath ID: 10400163}
\textbf{Source problem code:} catalogue entry
\subsection{Definitions and mathematical statement}
An extended topological quantum field theory (ETQFT) \(V\) has a category \(\mathcal C_V\) assigned to a circle. Write \(V_{\mathcal C}\) for the source's ETQFT construction from a modular tensor category \(\mathcal C\). T. Kerler asks for the following developments.
\begin{enumerate}
\item Determine how an ETQFT \(V\) with circle category \(\mathcal C\) can differ from \(V_{\mathcal C}\), and define a suitable equivalence of ETQFTs and modular tensor categories under which assignment of the circle category gives a bijective correspondence of equivalence classes.
\item Extend half-projectivity to the setting allowing surfaces that are both disconnected and have boundary.
\item Give constructions and axioms of ETQFTs under more relaxed boundedness or modularity requirements, including theories relevant to finite-type invariants.
\end{enumerate}
The accompanying source discussion is part of the specification: its functor \(\mathrm{Alg}\to\mathrm{Cob}_{\bullet}\) forces \(\mathcal C_V\) to satisfy approximately all modular tensor category axioms and to contain a Hopf algebra object with additional properties. Under an unspecified rigidity assumption this object must agree with the one used for \(V_{\mathcal C}\); different algebra structures on that object may still produce ambiguities, and reference [216] is said to exhibit several choices. Item (2) is to be informed by the double composition laws for surgery tangles in [216] and a generalization of [210].

For item (3), the source contrasts a non-semisimple quantum-group representation category with its canonical semisimple subquotient, referring to [208]: the resulting semisimple Reshetikhin--Turaev theory differs from the non-semisimple theory obtained by the Hennings algorithm. It also discusses starting with a rigid monoidal category \(\mathcal B\) without braiding, taking its Drinfel'd double \(D(\mathcal B)\) with suitable ribbon data and using \(V_{D(\mathcal B)}\), or, when \(\mathcal B\) is semisimple, constructing a Turaev--Viro theory \(W_{\mathcal B}\) from \(6j\)-symbols. The extract does not supply the hypotheses ensuring these constructions, the ETQFT domain and axioms, or the original half-projectivity law. Its extracted notation also fails to distinguish the two categories and theories in the semisimplification discussion. These contextual assertions therefore do not specify unrestricted construction theorems for arbitrary rigid monoidal categories.
\subsection{Short English statement}
Classify the possible extended theories with a fixed circle category, extend half-projectivity to disconnected surfaces with boundary, and broaden the axioms and constructions beyond the source's boundedness and modularity assumptions.
\subsection{Sources}
\begin{itemize}
\item T. Kerler, attributed Problem 8.17, in T. Ohtsuki, \emph{Problems on invariants of knots and 3-manifolds} (2002), supplied PDF page 132. \url{https://msp.org/gtm/2002/04/gtm-2002-04-024s.pdf}. Recorded access 2026-07-29; extraction pdf-font-structure.
\item The excerpt explicitly refers to [216] for surgery-tangle double composition laws and choices, [210] for the proposed half-projective generalization, and [208] for semisimple subquotients. Their bibliography entries are not supplied. It names Reshetikhin--Turaev, Hennings, Drinfel'd double and Turaev--Viro constructions without further complete references.
\item Local UnsolvedMath registry, \texttt{data/unsolvedmath/problems.json}, permanent ID \(10400163\), exact source code \texttt{catalogue entry}; supplied research classification OPEN-TRIAGE.
\end{itemize}
\end{document}
